\section{Proofs for Section 4}\label{appendix:proofs_conditional_composition}

In the following, we first confirm that the conditional composition~\cref{lemma:conditional_composition} immediately follows as a special case of the general conditional composition framework from~\cite{choquette2023privacy}. In the next sections, we then provide the omitted proofs for our instantiation of the conditional composition framework.

The general conditional composition result is:
\begin{theorem}[Theorem 3.1 in~\cite{choquette2023privacy}]\label{thm:conditioningtrick}
Let $\mathcal{M}_1: \mathcal{D} \rightarrow \mathcal{X}_1, \mathcal{M}_2: \mathcal{X}_1 \times \mathcal{D} \rightarrow \mathcal{X}_2, \mathcal{M}_3: \mathcal{X}_1 \times \mathcal{X}_2 \times \mathcal{D} \rightarrow \mathcal{X}_3, \ldots \mathcal{M}_N$ be a sequence of adaptive mechanisms, where each $\mathcal{M}_n$ takes a dataset in $\mathcal{D}$ and the output of mechanisms $\mathcal{M}_1, \ldots, \mathcal{M}_{n-1}$ as input. Let $\mathcal{M}$ be the mechanism that outputs $(x_1 = \mathcal{M}_1(D), x_2 = \mathcal{M}_2(x_1, D), \ldots, x_N = \mathcal{M}_N(x_1, \ldots, x_{N-1}, D))$. Fix any two adjacent datasets $D, D'$. 

Suppose there exists ``bad events'' $E_1 \subseteq \mathcal{X}_1, E_2 \subseteq \mathcal{X}_1 \times \mathcal{X}_2, \ldots... E_{N-1} \subseteq \mathcal{X}_1 \times \mathcal{X}_2 \times \ldots \times \mathcal{X}_{N-1}$ such that \[\Pr_{x \sim \mathcal{M}(D)}\left[\exists n: (x_1, x_2, \ldots x_n) \in E_n\right] \leq \delta\] and pairs of distributions $(P_1, Q_1), (P_2, Q_2), \ldots (P_N, Q_N)$ such that the PLD of $\mathcal{M}_1(D)$ and $\mathcal{M}_1(D')$ is dominated by the PLD of $P^{(1)}, Q^{(1)}$ and for any $n \geq 1$ and ``good'' output $(x_1, x_2, \ldots x_n) \notin E_n$, the PLD of $\mathcal{M}_{n+1}(x_1, \ldots, x_n, D)$ and $\mathcal{M}_{n+1}(x_1, \ldots, x_n, D')$ is dominated by the PLD of $P^{(n+1)}, Q^{(n+1)}$. Then for all $\epsilon$:

\[H_\epsilon(\mathcal{M}(D), \mathcal{M}(D')) \leq H_\epsilon\left(P^{(1)} \times P^{(2)} \times \ldots \times P^{(N)}, Q^{(1)} \times Q^{(2)} \times \ldots \times Q^{(N)}\right) + \delta.\]
\end{theorem}
In our case, we do not have a mechanism, but only its dominating pair $(P, Q)$ from~\cref{lem:dominating_pair}.
However, for any $D, D'$, we can define the dummy mechanism $\mathcal{M}(D) = P$ and $\mathcal{D}' = Q$ and $M_{n}(x_1,\dots,x_n, D) = P_{x_n \mid x_{1:n-1}}$ and $M_{n}(x_1,\dots,x_n, D') = Q_{x_n \mid x_{1:n-1}}$. The special case below then follows immediately.
\conditionalcomposition*


\subsection{Proofs for Section 4.1}\label{appendix:proofs_cond_comp_algorithm}

\reversehazarddominance*
\begin{proof}
    Let $F_i = P(z \leq i)$ denote the CDF of $z$. By definition, $\lambda_i = P(z=i \mid z \leq i) = P(z=i)/F_i$.
    Furthermore, $F_{i} = F_{i+1} - P(z=i+1)$.
    It thus follows that $F_{i} = F_{i+1} - P(z=i+1) = F_{i+1}(1 - \lambda_{i+1})$. Since $F_i=1$ for any $i \geq b$, recursive application of this equality yields $F_i = \prod_{j=i+1}^b (1 - \lambda_j)$.
    The condition $\lambda_j \leq \lambda'_j$ implies $(1 - \lambda_j) \geq (1 - \lambda'_j)$ for all $j$. Consequently, $F_i \geq F'_i$ for all $i$, i.e., $p(z)$ is stochastically dominated by $p'(z)$.
\end{proof}

For the next result (\cref{corollary:reverse_hazard_to_dominance}), we first recall~\cref{lemma:stochastic_dominance_to_dp_dominance}, which relates stochastic dominance between mixture weights to dominance between Gaussian mixture mechanisms in the differential privacy sense:
\mogweightmonotonicity*

\reversehazardtodominance*
\begin{proof}
    Let $p(z)$ denote the conditional mixture weights $p(z \mid \vy_{1:n-1})$ and use $\lambda_i = P(z = i \mid i \leq i, \vy_{1:n-1})$ to denote the corresponding reverse hazard function.
    By assumption, the unconditional mixture weights we construct are $p^{(n)}(z = i) = \overline{\lambda_i} \prod_{j=i+1}^b (1 - \overline{\lambda_i})$.
    From the definition of conditional probability and the recurrence relation $F_i = \prod_{j=i+1}^b (1 - \lambda_j)$ used in our previous proof, it follows that the corresponding reverse hazard function is:
    \begin{equation*}
        P^{(n)}(z = i \mid z \leq i) = \frac{p^{(n)}(z = i)}{F_i} =
            \frac{
                \overline{\lambda_i} \prod_{j=i+1}^b (1 - \overline{\lambda_i})
            }{
                \prod_{j=i+1}^b (1 - \lambda_j)
            }
            = \overline{\lambda_i}.
    \end{equation*}
    That is, $\overline{\lambda_i}$ is exactly the reverse hazard function of $p^{(n)}$.
    By assumption, $\lambda_i \leq \overline{\lambda_i}$.
    It thus follows from~\cref{lemma:reverse_hazard_to_stochastic_dominance}
    that $p(z)$ is stochastically dominated by $p^{(n)}$.
    The result is then an immediate consequence of~\cref{lemma:stochastic_dominance_to_dp_dominance}.
\end{proof}

\reversehazardloss*
\begin{proof}
    Let $p(\vy_{1:n-1} \mid z=j) = \mathcal{N}(\vy_{1:n-1} \mid \bm{m}_{j,:n-1}, \sigma^2 \eye)$ denote the likelihood of the earlier outcome $\vy_{1:n-1}$ conditioned on mixture component $j$.
    Let $p(z=j) = \frac{1}{b}$ denote the uniform prior that arises from random allocation with $b$ batches per epoch.
    Using the definition of conditional probability and Bayes' rule, the conditional reverse hazard function is:
    \begin{align*}
        &\Pr(z=i \mid z \leq i, \vy_{1:n-1})\\
        = &\frac{p(\vy_{1:n-1} \mid z=i)  p(z=i)}{\sum_{j=1}^i p(\vy_{1:n-1} \mid z=j)  p(z=j)}\\
        =& \frac{p(\vy_{1:n-1} \mid z=i)}{\sum_{j=1}^i p(\vy_{1:n-1} \mid z=j)}\\
        = & \left(1 + \frac{\sum_{j=1}^{i-1} p(\vy_{1:n-1} \mid z=j)}{p(\vy_{1:n-1} \mid z=i)}\right)^{-1}.
    \end{align*}
    Recall that $L_i(\vy_{1:n-1}) = \log \left( \frac{\frac{1}{i-1}\sum_{j=1}^{i-1} p(\vy_{1:n-1} \mid z=j)}{p(\vy_{1:n-1} \mid z=i)} \right)$.
    Multiplying and dividing the second summand by $i-1$ yields:
    \begin{equation}
        \frac{\sum_{j=1}^{i-1} p(\vy_{1:n-1} \mid z=j)}{p(\vy_{1:n-1} \mid z=i)}
        = (i-1) \exp(L_i(\vy_{1:n-1}))
        = \exp(\log(i-1) + L_i(\vy_{1:n-1})).
    \end{equation}
    Applying the definition of the sigmoid function $s(x) = (1+e^{-x})^{-1}$, we obtain:
    \begin{equation}
        \Pr(z=i \mid z \leq i, \vy_{1:n-1}) = s\left(-\big(\log(i-1) + L_i(\vy_{1:n-1})\big)\right),
    \end{equation}
    which completes the proof.
\end{proof}

While~\cref{section:per_step_dominating_pairs} in its entirety already served as a proof sketch,
we now provide a formal proof for the correctness of~\cref{algorithm:conditional_composition}:
\condcompbothcorrect*
\begin{proof}
    \cref{algorithm:conditional_composition} allocates a failure probability $\beta_i$ to each component $i$ such that $\sum_{i} \beta_i = \beta$.
    For each $i$, the subroutine \texttt{whp\_lower} returns a threshold $\tau_i$ such that $\Pr[L_i(\vy_{1:n-1}) < \tau_i] \leq \beta_i$, where the probability is over the randomness of the privacy loss random variable $L_{\tilde{P},\tilde{Q},\tilde{R}}$ generated by the algorithm's reference distribution $\tilde{R}$.
    In the ``remove'' case, $\tilde{R}$ is the joint distribution $P_{\vy_{1:n-1}}$ of the first $n-1$ outcomes of $\vy \sim P$.
    In the ``add'' case, $\tilde{R}$ is the joint distribution $Q_{\vy_{1:n-1}}$ of the first $n-1$ outcomes of $\vy \sim Q$.


    Define the failure event $\overline{A} = \bigcup_i \{ L_i(\vy_{1:n-1}) < \tau_i \}$
    with $L_i$ defined in~\cref{lemma:reverse_hazard_loss}. By a union bound, $\Pr(\overline{A}) \leq \sum \beta_i = \beta$.
    
    Conditioned on the success event $A = \overline{\overline{A}}$, we have $L_i(\vy_{1:n-1}) \geq \tau_i$ for all $i$.
    Consequently, by the monotonicity of the sigmoid function and the derivation in \cref{lemma:reverse_hazard_loss}, the true conditional reverse hazard satisfies:
    \begin{equation}
        \Pr(z=i \mid z \leq i, \vy_{1:n-1}) = s(-\log(i-1) - L_i) \leq s(-\log(i-1) - \tau_i) = \overline{\lambda_i}.
    \end{equation}
    Since the means are sorted non-decreasingly (Line 3 of~\cref{algorithm:conditional_composition}), the conditions of \cref{corollary:reverse_hazard_to_dominance} (larger reverse hazard implies dominance in the differential privacy sense) are satisfied. Thus, in the ``remove'' case the returned pair $(P^{(n)}, Q^{(n)})$ dominates the conditional distributions $P_{y_n | \tilde{y}_{1:n-1}}$
    and $Q_{y_n | \tilde{y}_{1:n-1}}$ with probability at least $1-\beta$ under $P$.
    In the ``add'' case, the returned pair $(Q^{(n)}, P^{(n)})$ dominates the conditional distributions  $Q_{y_n | \tilde{y}_{1:n-1}}$
    and $P_{y_n | \tilde{y}_{1:n-1}}$ with probability at least $1-\beta$ under $Q$.
\end{proof}

\subsection{Proofs for Section 4.2}\label{appendix:proofs_variational}
We conclude by proving our variational bounds on the privacy loss, including the previously omitted ``remove'' direction.
\variationalgeneric*
\begin{proof}
    Let $p$ be the density of $P$ and $q$ be the density of $Q$.
    By definition of the privacy loss $L_i$ and the product rule for logarithms, we have
    \begin{equation*}
        L_i(\vx) = \log \left( \frac{p(\vx)}{q(\vx)} \right) =
        \log \left( \sum_{j=1}^{i-1} \mathcal{N}(\vx \mid \vmu_j, \sigma^2 \eye) \pi_i \right)
        - \log(\mathcal{N}(\vx \mid \vmu_i, \sigma^2 \eye).
    \end{equation*}
    The first term is simply the log-likelihood of a Gaussian mixture model.
    It is, quasi by definition, lower-bounded by the evidence lower bound (for the formal proof, see, e.g., Section 2.2 of~\cite{blei2017variational}).
    We thus have
    \begin{align*}
        L_i(\vx)
        & \geq \mathbb{E}_{j \sim \psi}\left[\log \left( \mathcal{N}(\vx \mid \vmu_j, \sigma^2 \eye) \right)\right] - \mathrm{KL}( \psi || \pi) - \log(\mathcal{N}(\vx \mid \vmu_i, \sigma^2 \eye)
        \\
        &
        =
         \mathbb{E}_{j \sim \psi}\left[\log \left( \mathcal{N}(\vx \mid \vmu_j, \sigma^2 \eye) \right)  - \log(\mathcal{N}(\vx \mid \vmu_i, \sigma^2 \eye) \right] - \mathrm{KL}( \psi || \pi)
        \\
        &
        =
         \mathbb{E}_{j \sim \psi}\left[ L_{i,j}(\vx)\right] - \mathrm{KL}( \psi || \pi)
         \\
         &
         = \underline{L_i}(\vx),
    \end{align*}
    where the first equality is due to linearity of expectation.
\end{proof}

Next, we use this bound on the privacy loss $L_i$ to derive a lower tail bound on $L_i(\vx)$ with $\vx \sim \tilde{R}$ where $\tilde{R}$ is an arbitrary Gaussian mixture.
The concrete instantiations for the ``remove'' and ``add'' $\tilde{R}$ from~\cref{algorithm:conditional_composition} will then follow as special cases.
\begin{lemma}\label{lemma:variational_distribution_general}
    Let $\vx \sim \tilde{R}$ with $\tilde{R} = \sum_{k=1}^K \mathcal{N}(\vrho_k, \sigma^2 \eye) \frac{1}{K}$.
    Let $\Psi = \{\psi^{(1)}, \dots, \psi^{(K)}\}$ be a collection of variational distributions on $[i-1]$.
    Consider the random variable $Y$ generated by sampling a component index $k \sim \mathrm{Uniform}([K])$, sampling $\vx \sim \mathcal{N}(\vrho_k, \sigma^2 \eye)$, and computing the component-specific lower bound $\underline{L}_i(\vx; \psi^{(k)})$ as defined in~\cref{eq:variational_bound}.
    Then $Y \sim \sum_{k=1}^K \mathcal{N}(\nu_k, \xi_k) \frac{1}{K}$ with component-specific means $\nu_k \in \sR$
    and standard deviations $\xi_k \in \sR_+$ defined by:
    \begin{align*}
        \nu_k &= \frac{1}{\sigma^2} \vrho_k^T (\mathbb{E}_{j \sim \psi^{(k)}}[\vmu_j] - \vmu_i) + \frac{1}{2\sigma^2} \left( \|\vmu_i\|_2^2 - \mathbb{E}_{j \sim \psi^{(k)}}[\|\vmu_j\|_2^2] \right) - \mathrm{KL}(\psi^{(k)} \mid\mid \pi), \\
        \xi_k &= \frac{1}{\sigma} \| \vmu_i - \mathbb{E}_{j \sim \psi^{(k)}}[\vmu_j] \|_2.
    \end{align*}
\end{lemma}
\begin{proof}
    We first derive a bound for each component $k$. Recall that the per-component privacy loss is $L_{i,j}(\vx) = \frac{1}{\sigma^2} \vx^T (\vmu_j - \vmu_i) + \frac{1}{2\sigma^2} \left( \|\vmu_i\|_2^2 - \|\vmu_j\|_2^2 \right)$.
    For a specific mixture component $k$, we use the variational distribution $\psi^{(k)}$. By linearity of expectation, the lower bound $\underline{L}_i(\vx; \psi^{(k)}) = \mathbb{E}_{j \sim \psi^{(k)}} [L_{i,j}(\vx)] - \mathrm{KL}(\psi^{(k)} \mid\mid \pi)$ from~\cref{lemma:variational_generic} can be expressed as an affine transformation:
    \begin{equation*}
        \underline{L}_i(\vx; \psi^{(k)}) = \vx^T \underbrace{\left[ \frac{1}{\sigma^2} (\mathbb{E}_{j \sim \psi^{(k)}}[\vmu_j] - \vmu_i) \right]}_{\mathbf{a}_k} + \underbrace{\frac{1}{2\sigma^2} \left( \|\vmu_i\|_2^2 - \mathbb{E}_{j \sim \psi^{(k)}}[\|\vmu_j\|_2^2] \right) - \mathrm{KL}(\psi^{(k)} \mid\mid \pi)}_{c_k}.
    \end{equation*}
    Conditioned on the latent component $z=k$, the vector $\vx$ follows $\mathcal{N}(\vrho_k, \sigma^2 \eye)$. Thus, the conditional distribution of the lower bound is univariate Gaussian with mean $\nu_k$ and variance $\xi_k^2$.
    The mean can be calculated via:
    \begin{align*}
        \nu_k &= \mathbb{E}_{\vx \sim \mathcal{N}(\vrho_k, \sigma^2 \eye)} [\vx^T \mathbf{a}_k + c_k] = \vrho_k^T \mathbf{a}_k + c_k \\
        &= \frac{1}{\sigma^2} \vrho_k^T (\mathbb{E}_{j \sim \psi^{(k)}}[\vmu_j] - \vmu_i) + \frac{1}{2\sigma^2} \left( \|\vmu_i\|_2^2 - \mathbb{E}_{j \sim \psi^{(k)}}[\|\vmu_j\|_2^2] \right) - \mathrm{KL}(\psi^{(k)} \mid\mid \pi).
    \end{align*}
    The variance $\xi_k^2$ can be calculated via:
    \begin{align*}
        \xi_k^2 &= \mathrm{Var}(\vx^T \mathbf{a}_k + c_k \mid \vx \sim \mathcal{N}(\vrho_k, \sigma^2 \eye)) = \mathbf{a}_k^T (\sigma^2 \eye) \mathbf{a}_k \\
        &= \sigma^2 \left\| \frac{1}{\sigma^2} (\mathbb{E}_{j \sim \psi^{(k)}}[\vmu_j] - \vmu_i) \right\|_2^2 \\
        &= \frac{1}{\sigma^2} \| \mathbb{E}_{j \sim \psi^{(k)}}[\vmu_j] - \vmu_i \|_2^2.
    \end{align*}
    Taking the square root yields $\xi_k = \frac{1}{\sigma} \| \vmu_i - \mathbb{E}_{j \sim \psi^{(k)}}[\vmu_j] \|_2$. Since the mixture component $k$ is chosen uniformly at random with probability $1/K$, the marginal distribution is the mixture $\frac{1}{K} \sum_{k=1}^K \mathcal{N}(\nu_k, \xi_k^2)$.
\end{proof}
Applying this result to the ``remove'' case, we have:
\begin{theorem}\label{thm:remove_tail_bound}
    Let $\vx \sim \tilde{R}$ with $\tilde{R} = \frac{1}{b} \sum_{k=1}^b \mathcal{N}(\vmu_k, \sigma^2 \eye)$.
    Let $\Psi = \{\psi^{(1)}, \dots, \psi^{(b)}\}$ be a collection of variational distributions on $[i-1]$.
    For any threshold $\tau \in \sR$, the left-tail probability of the privacy loss random variable 
    $L_i(\vx)$ with $L_i$ defined in~\cref{lemma:reverse_hazard_loss} is bounded by:
    \begin{equation}
        \Pr_{\vx \sim \tilde{R}}[L_i(\vx) \leq \tau] \leq \frac{1}{b} \sum_{k=1}^b \Phi\left( \frac{\tau - \nu_k}{\xi_k} \right),
    \end{equation}
    where $\Phi$ is the cumulative distribution function of the standard normal distribution, and per-component parameters are:
    \begin{align*}
        \nu_k &= \frac{1}{\sigma^2} \vmu_k^T (\mathbb{E}_{j \sim \psi^{(k)}}[\vmu_j] - \vmu_i) + \frac{1}{2\sigma^2} \left( \|\vmu_i\|_2^2 - \mathbb{E}_{j \sim \psi^{(k)}}[\|\vmu_j\|_2^2] \right) - \mathrm{KL}(\psi^{(k)} \mid\mid \pi), \\
        \xi_k &= \frac{1}{\sigma} \| \vmu_i - \mathbb{E}_{j \sim \psi^{(k)}}[\vmu_j] \|_2.
    \end{align*}
\end{theorem}
\begin{proof}
    By the Law of Total Probability, $\Pr_{\vx \sim \tilde{R}}[L_i(\vx) \leq \tau] = \frac{1}{b} \sum_{k=1}^b \Pr_{\vx \sim \mathcal{N}(\vmu_k, \sigma^2 \eye)} [L_i(\vx) \leq \tau]$.
    Since $L_i(\vx) \geq \underline{L}_i(\vx; \psi^{(k)})$, we have the bound
    $\Pr[L_i(\vx) \leq \tau] \leq \Pr[\underline{L}_i(\vx; \psi^{(k)}) \leq \tau]$ for each component $k$.
    Applying \cref{lemma:variational_distribution_general} with $\vrho_k = \vmu_k$, the variable $\underline{L}_i(\vx; \psi^{(k)})$ conditioned on component $k$ follows $\mathcal{N}(\nu_k, \xi_k)$.
    Consequently, each term in the sum is exactly $\Phi\left( \frac{\tau - \nu_k}{\xi_k} \right)$.
\end{proof}
For the ``add'' case, we observe that the lower bound on our privacy loss is a single, univariate Gaussian:
\variationalboundadd*
\begin{proof}
    This immediately follows as a special case of~\cref{lemma:variational_distribution_general} with $\vrho_k = \vzero$ and $K=1$.
\end{proof}
Naturally, this result directly implies the following analytical tail bound for the ``add'' case:
\begin{corollary}\label{thm:add_tail_bound}
    Let $\vx \sim \tilde{R}$ with $\tilde{R} = \mathcal{N}(\vzero, \sigma^2 \eye)$.
    Let $\psi$ be an  arbitrary variational distribution on $[i-1]$.
    For any threshold $\tau \in \sR$, the left-tail probability of the privacy loss random variable 
    $L_i(\vx)$ with $L_i$ defined in~\cref{lemma:reverse_hazard_loss} is bounded by:
    \begin{equation}
        \Pr_{\vx \sim \tilde{R}}[L_i(\vx) \leq \tau] \leq \Phi\left( \frac{\tau - \nu}{\xi} \right),
    \end{equation}
    where $\Phi$ is the cumulative distribution function of the standard normal distribution, with parameters
    \begin{align*}
        \nu &= \frac{1}{2\sigma^2} \left( \|\vmu_i\|_2^2 - \mathbb{E}_{j \sim \psi}[\|\vmu_j\|_2^2] \right) - \mathrm{KL}(\psi \mid\mid \pi), \\
        \xi &= \frac{1}{\sigma} \| \vmu_i - \mathbb{E}_{j \sim \psi}[\vmu_j] \|_2.
    \end{align*}
\end{corollary}
\begin{remark}[Handling Identical Components]
    In structured matrix mechanisms, multiple batches often share the same mixture mean, i.e., $\vmu_j = \vmu_i$. For these components, the privacy loss $L_{i,j}(\vx)$ is identically zero for all $\vx$. Including them in the variational mean $\nu$ is suboptimal as it pulls the lower bound toward zero. We handle this by setting the variational weights $\psi_j = 0$ for all such indices. This effectively pulls the zero-loss components out of the expectation and into the KL divergence term, where they contribute a constant offset to the mean $\nu$.
\end{remark}